\documentclass[11pt]{article}
\usepackage[margin=1in]{geometry}
\usepackage{amsmath,amssymb,amsthm,mathtools}
\usepackage[hidelinks]{hyperref}
\newtheorem{theorem}{Theorem}
\newtheorem{lemma}[theorem]{Lemma}
\newtheorem{corollary}[theorem]{Corollary}
\newcommand{\Q}{\mathbb Q}
\newcommand{\C}{\mathbb C}
\newcommand{\Lfac}{\mathcal L}
\title{Consecutive moment windows, polynomial markers, and the Rocks sources}
\author{Internal source audit and exact implementations; no priority claim}
\date{16 November 2026}
\begin{document}
\maketitle
\noindent\textbf{Scope.}
The earlier notes prove absence of an infinite zero tail for several
classes. Here finite consecutive windows are proved for powers of linear
forms, with explicit integer-Gamma and polynomial-marker extensions.
The exact rational generating function decides whether a marked sequence
is identically zero and otherwise gives a certified window at every index.
This does not prove unrestricted Strong Factorial.

The principal credit is Brady Rocks. The author-hosted
\href{sources/rocks-2015-author-preprint.pdf}{September 2015, 28-page preprint}
is \emph{not} his dissertation; the latter has its own
\href{https://doi.org/10.7936/K7BR8QB4}{Washington University record}.
The journal article is \emph{J. Algebra} 478 (2017), 341--352,
\href{https://doi.org/10.1016/j.jalgebra.2017.01.039}{DOI 2017.01.039}.
Its published abstract confirms the linear-power, prime-power and
construction results. The full final publisher text was unavailable;
the proof audit below is of the pinned author preprint, with its numbering
and explicitly corrected affine-window statement. No identity between
draft and final text, priority claim, or outside review is asserted.

\section{Reciprocal sections cannot have a long zero window}
Let $K$ be any field and $P(z)\in K[z]$ have $P(0)=1$ and degree $N>0$.
Over an algebraic closure write $P(z)=\prod_{i=1}^N(1-c_i z)$,
where every $c_i\ne0$ and repetitions are allowed. For an integer $r\ge1$ put
\[
 Q(t)=\prod_{i=1}^N(1-c_i^r t),\qquad
 W(z)=\frac{Q(z^r)}{P(z)}
     =\prod_{i=1}^N\sum_{j=0}^{r-1}(c_i z)^j.
\]
Write $P(z)^{-1}=\sum_{k\ge0}a_kz^k$ and $W(z)=\sum_jw_jz^j$.
Extracting exponents divisible by $r$ gives
\[
 \sum_{m\ge0}a_{rm}t^m=
 \frac{\sum_{m\ge0}w_{rm}t^m}{Q(t)}.                 \tag{1}
\]
The numerator has constant term one and degree at most
$\lfloor(r-1)N/r\rfloor<N$. Thus (1) is nonzero and proper.

If a rational series $B(t)=R(t)/D(t)$ is proper, $D(0)=1$, and
$\deg D=d$, its coefficients satisfy the homogeneous order-$d$ recurrence
from $D(t)B(t)=R(t)$ at every index at least $d$. A zero block of length
$d$, starting at an index $n\ge1$, propagates to an infinite zero tail.
Then $B$ is polynomial; a proper rational function can be polynomial only
if it is zero. This proves Rocks's Proposition 3.4/Corollary 3.6 in a
form retaining an executable rational identity:

\begin{theorem}\label{thm:sections}
For every $n,r\ge1$, some coefficient among
$a_{rn},a_{r(n+1)},\ldots,a_{r(n+N-1)}$ is nonzero.
The argument works in every characteristic and does not divide by $r$.
\end{theorem}

This is a multiplicative inverse statement. It is distinct from Furter's
compositional inverse theorem, although both give missing-coefficient bounds.

\section{Gamma windows and every polynomial marker}
Let $Y_i$ be independent $\operatorname{Gamma}(\alpha_i,1)$ variables,
where the $\alpha_i$ are positive integers, and put $L=\sum_i c_iY_i$.
Assume some $c_i\ne0$. Formal Taylor expansion, justified coefficientwise
by the elementary Gamma moments, gives
\[
 \sum_{k\ge0}\frac{\mathbb E[L^k]}{k!}z^k
 =\prod_i(1-c_i z)^{-\alpha_i}=P(z)^{-1}.             \tag{2}
\]
Its actual polynomial denominator degree is
$N=\sum_{c_i\ne0}\alpha_i$; equal nonzero nodes combine multiplicities.
For $F=L^r$, Theorem~\ref{thm:sections} proves that every $N$ consecutive
positive moments contain a nonzero one. With all $\alpha_i=1$, $N$ is the
number of active variables, at most the monomial count of $L^r$.
Consequently this recovers Rocks's Strong Factorial theorem for every
power of a linear form, not merely ordinary FC.

More generally individual shapes can be arbitrary positive numbers if their
sum at each distinct nonzero node $c$ is an integer. Grouping equal nodes
in (2) gives the same polynomial $P$. A monomial marker shifts those group
sums by integers, so all the conclusions below still apply. The exporter
supports exact positive rational individual shapes with this condition;
in particular half-integer shapes may aggregate to integer blocks. No
integrality condition is needed at node zero. This is not an extension to
arbitrary noninteger shape sums at distinct active nodes.

Let $M(Y)=\sum_b m_bY^b$ be any polynomial marker. Its ordinary exponential
moment series is the explicit rational function
\[
 K_M(z)=\sum_{k\ge0}\frac{\mathbb E[M(Y)L^k]}{k!}z^k
 =\sum_b m_b\prod_i(\alpha_i)_{b_i}
                       (1-c_i z)^{-\alpha_i-b_i}.    \tag{3}
\]
Every summand is proper because $L\ne0$. Thus $K_M$ is proper or zero.
A common denominator is
$P_M(z)=\prod_{c_i\ne0}(1-c_i z)^{\alpha_i+\deg_{Y_i}M}$.
Repeated nodes and cancellations may reduce it.

For any proper $K_M=R/P_M$, the same section operation as (1) gives
\[
 B_{M,r}(t)=\sum_{m\ge0}
       \frac{\mathbb E[M(Y)L^{rm}]}{(rm)!}t^m
       =\frac{\operatorname{section}_r
          \{R(z)Q_M(z^r)/P_M(z)\}}{Q_M(t)},           \tag{4}
\]
where $Q_M(t)=\prod_{c_i\ne0}(1-c_i^r t)^{\alpha_i+\deg_{Y_i}M}$.
The polynomial inside the section has degree strictly below
$r\deg P_M$, so (4) is proper, including after cancellation. This proves:

\begin{theorem}\label{thm:marked}
For every marker $M$, exactly one of the following holds. Either every
marked moment in (4), including $m=0$, is zero, or the reduced rational
function has denominator degree $d\ge1$ and every $d$ consecutive positive
marked moments contain a nonzero one. One valid upper bound is
$\sum_{c_i\ne0}(\alpha_i+\deg_{Y_i}M)$.
The reduced denominator degree is a certified window, not a claim that
it equals the smallest possible zero-gap bound.
\end{theorem}

The zero branch genuinely occurs. With shapes $(1,1)$, nodes $(1,-1)$,
$r=2$ and marker $M=Y_1-Y_2$, odd symmetry makes every selected marked
moment zero, even though the unit moments of $L^2$ are nonzero.
The exporter proves this by a zero numerator in (4). An identically zero
marked sequence is not a pure-null phase.

On a Dirichlet simplex, a polynomial marker can first be homogenized on
$\sum T_i=1$ to degree $e$. For its lift $\widetilde M(Y)$,
\[
 \mathbb E_\Gamma[\widetilde M(Y)L^{rm}]
  =(A)_{e+rm}\mathbb E_{\rm Dir}[M(T)\ell(T)^{rm}].
\]
The radial factor never vanishes, so the same certified windows transfer.
The Gaussian substitution and orthogonal marker averaging from the earlier
bridge note retain the corresponding original Gaussian markers; monomials
with an odd coordinate average to zero and even ones give polynomial Gamma
markers. These are structured phase results, not general Gaussian MZ.

\section{Affine shifts and new independent variables}
For a nonzero constant $a$, let $f_k=\mathbb E[(a+L)^k]/k!$.
Equation (2) becomes
\[
 P(z)\sum_{k\ge0}f_kz^k=e^{az},\qquad
 \sum_{j=0}^Np_jf_{k-j}=a^k/k!\ne0\quad(k\ge N).
\]
An $N+1$ consecutive zero block is impossible. For factorial shapes this
is the correct $m+1$-term affine SFC window when $m$ variable terms and a
nonzero constant are present.

\textbf{Preprint correction.} The displayed statement of Theorem 3.7 on
preprint page 9 says $m$ moments. Its proof uses $m+1$ and proves the
correct SFC assertion. The shorter displayed assertion is false: $Y-1$
is nonzero and has first factorial moment zero. The page image was checked;
this is not an extraction error. We do not attribute the draft typo to
the unavailable final journal text.

Rocks's Theorem 3.12 also has a useful precise form. Suppose a nonzero
phase $F$ has a known window bound $B$ for its factorial moments. For a
new independent exponential variable $U$ and $G=F+\lambda U$, set
$f_k=\mathbb E[F^k]/k!$, $g_k=\mathbb E[G^k]/k!$. Independence gives
$g_k=\sum_{j=0}^k\lambda^jf_{k-j}$, hence
$f_k=g_k-\lambda g_{k-1}$. A length-$B+1$ zero block for $g$ would give
a length-$B$ zero block for $f$. Thus $G$ has bound $B+1$; for a new
nonzero linear term its monomial count also increases by one, preserving
SFC. More generally a new integer-shape Gamma variable of shape $a$
contributes $(1-\lambda z)^{-a}$ to this normalized moment series.
Applying the degree-$a$ difference polynomial gives the bound $B+a$.
Zero added coefficients may simply be omitted.

\section{The common-prime power theorem: checked quantifiers}
Fix one prime $p$ and positive integers $\beta_i$, and set
$G=\sum_{i=1}^m\lambda_iY_i^{p^{\beta_i}}$ for factorial variables.
Rocks's Theorem 3.10 proves that $G^r$ has window $m$ for every $r\ge1$.
The exponents are powers of the \emph{same} prime. This does not assert
the mixed-prime case such as $Y_1^2+\lambda Y_2^3$.

Here is the checked all-index argument. Write
$b_i=(p^{\beta_i}-1)/(p-1)$. Splitting a factorial into multiples of $p$
and the other factors, and iterating, gives
\[
 \frac{(p^{\beta_i}k)!}{p^{b_i k}k!}\in\mathbb Z,
 \qquad
 \frac{(p^{\beta_i}k)!}{p^{b_i k}k!}
       \equiv(-1)^{b_i k}\pmod p.                  \tag{5}
\]
The last congruence is Wilson's theorem applied in each block. Expand
$\mathbb E[G^n]/n!$ as a homogeneous polynomial in the coefficients and
substitute $\lambda_i=(-p)^{-b_i}x_i$. Its coefficients are integers by
(5), and its reduction is the complete homogeneous polynomial $h_n(x)$.
The sampled window $n,\ldots,n+m-1$ for $G^r$ therefore reduces to
$h_{rn},\ldots,h_{r(n+m-1)}$. Theorem~\ref{thm:sections}, over
$\overline{\mathbb F}_p$, says their only common affine zero is the origin.

For completeness the homogeneous lifting is legitimate: in each sufficiently
large degree, the quotient by these homogeneous equations is a finitely
generated $\mathbb Z_{(p)}$-module whose reduction is zero. Nakayama makes
that degree quotient zero over $\mathbb Z_{(p)}$, hence over $\Q$ and $\C$.
Equivalently the proper projective zero scheme has empty special fibre and
therefore empty generic fibre. The use of the algebraic closure of the
residue field is essential; preprint Lemma 3.9 does contain that bar,
confirmed on its page image. Testing just $\mathbb F_p$-rational points
would not suffice.

The preprint's Proposition 3.11 treats an added constant only for starting
indices $n+m<p$. It is a finite-index-range statement, not full affine
prime-power SFC. Further two-monomial resultant and Newton-polygon results
in Sections 3.4 and following have not all been audited here.

\section{Filaseta--Rocks and the remaining two-term problem}
The publisher confirms Michael Filaseta and Brady Rocks,
\emph{Colloq. Math.} 145 (2016), 307--314,
\href{https://doi.org/10.4064/cm6761-3-2016}{DOI cm6761-3-2016}.
Their \href{sources/filaseta-rocks-author-preprint.pdf}{author preprint}
studies $f_{n,d}(x)=\sum_{j=0}^n\binom nj(dj)!x^j$ for fixed $d$,
and proves a lower density at least $\log2$ of irreducible indices $n$.
This is a density in $n$ for each fixed $d$, not all indices, all $d$
uniformly, or a complete two-term SFC proof.

The precise moment connection can be checked independently:
\[
 f_{n,d}(\lambda)=\mathbb E[(1+\lambda Y^d)^n],\qquad
 n f_{n-1,d}=n f_{n,d}-x f'_{n,d}.                 \tag{6}
\]
Since $f_{n,d}(0)=1$, adjacent polynomials have a common root exactly when
$f_{n,d}$ has a multiple root. Irreducibility in characteristic zero implies
squarefreeness. Thus the published density result certifies the
two-consecutive-moment assertion at a set of terminal indices of lower
density at least $\log2$, for each fixed $d$. It does not settle every
window. The source's number-theoretic density proof is mapped here, not
claimed as independently reconstructed or used in our new window exporter.

\section{Verification and continuing boundary}
\href{moment_windows.py}{The exact exporter} produces (4), the cancelled
denominator and a zero-sequence or finite-window certificate for algebraic
input coefficients. It retains all actual original marker variables.
Controls compare these identities with direct factorial/Gamma/simplex
moments, check repeated and zero nodes, sample aliasing, the affine draft
counterexample, common-prime factorial reductions, independent-variable
differences and (6). Finite controls do not replace the all-index proofs.

The preceding cubic two-coordinate proof was revalidated against every
recorded hash. A separate exact probe found common roots for 12 of the
125 exponent residues modulo five, despite invertible fifth-power
multiplication matrices. Thus the tempting blanket all-residue mod-five
upgrade fails. No general exceptional-scale theorem is inferred from the
125 rational gcds being one in the tested box; the written arithmetic
boundary remains unchanged.

General FC in at least three variables, unrestricted SFC, full two-monomial
SFC, all remaining EIT V/VI inputs and the complete bibliography remain
open. This note records progress on genuine finite windows and corrects
source identification; it does not redefine the full goal.
\end{document}
